%% ****** Start of file apstemplate.tex ****** %
%%
%%
%%   This file is part of the APS files in the REVTeX 4.2 distribution.
%%   Version 4.2a of REVTeX, January, 2015
%%
%%
%%   Copyright (c) 2015 The American Physical Society.
%%
%%   See the REVTeX 4 README file for restrictions and more information.
%%
%
% This is a template for producing manuscripts for use with REVTEX 4.2
% Copy this file to another name and then work on that file.
% That way, you always have this original template file to use.
%
% Group addresses by affiliation; use superscriptaddress for long
% author lists, or if there are many overlapping affiliations.
% For Phys. Rev. appearance, change preprint to twocolumn.
% Choose pra, prb, prc, prd, pre, prl, prstab, prstper, or rmp for journal
%  Add 'draft' option to mark overfull boxes with black boxes
%  Add 'showkeys' option to make keywords appear
%\documentclass[aps,prl,preprint,groupedaddress]{revtex4-2}
\documentclass[aps,prl,twocolumn,groupedaddress]{revtex4-2}
%\documentclass[aps,prl,preprint,superscriptaddress]{revtex4-2}
%\documentclass[aps,prl,reprint,groupedaddress]{revtex4-2}

% You should use BibTeX and apsrev.bst for references
% Choosing a journal automatically selects the correct APS
% BibTeX style file (bst file), so only uncomment the line
% below if necessary.
%\bibliographystyle{apsrev4-2}

\usepackage{graphicx}
\usepackage{epstopdf}
\usepackage{caption}
%\usepackage{amsmath}% http://ctan.org/pkg/amsmath
%\usepackage{amsthm}
%\usepackage{amsfonts}
%\usepackage{subfigure}
%\usepackage{hhline}
%\usepackage[miktex]{gnuplottex}
%\usepackage{xcolor}
\usepackage{amssymb}
\usepackage{amsmath}
\usepackage{color}
\usepackage{hyperref}
%\usepackage[percent]{overpic}
\usepackage{tikz}
\usepackage{mathrsfs}
\usepackage{wasysym}
\usepackage{tikz-cd}
%\usepackage{linearb}
%\usepackage{wsuipa}
\usepackage{calc}
\usepackage{tipa}

%\usepackage{stix} %\fisheye
\usepackage{stackengine,scalerel}

% so sections, subsections, etc. become numerated.
\setcounter{secnumdepth}{3}

\newcommand{\FF}{\mathbb{F}}
\newcommand{\NN}{\mathbb{N}}
\newcommand{\ZZ}{\mathbb{Z}}
\newcommand{\QQ}{\mathbb{Q}}
\newcommand{\RR}{\mathbb{R}}
\newcommand{\CC}{\mathbb{C}}
\newcommand{\II}{\mathbb{I}}
\newcommand{\GG}{\mathbb{G}}

\DeclareMathOperator*{\argmax}{arg\,max}
\DeclareMathOperator*{\argmin}{arg\,min}
\newcommand{\avrg}[1]{\left\langle #1 \right\rangle}
\newcommand{\nelta}{\bar{\delta}}
\newcommand{\bra}[1]{\left\langle #1\right|}
\newcommand{\ket}[1]{\left| #1 \right\rangle}
\newcommand{\sbra}[1]{\langle #1|}
\newcommand{\sket}[1]{| #1 \rangle}
\newcommand{\bek}[3]{\left\langle #1 \right| #2 \left| #3 \right\rangle}
\newcommand{\sbek}[3]{\langle #1 | #2 | #3 \rangle}
\newcommand{\braket}[2]{\left\langle #1 \middle| #2 \right\rangle}
\newcommand{\ketbra}[2]{\left| #1 \middle\rangle \middle\langle #2  \right|}
\newcommand{\sbraket}[2]{\langle #1 | #2 \rangle}
\newcommand{\sketbra}[2]{| #1 \rangle  \langle #2 |}
\newcommand{\norm}[1]{\left\lVert#1\right\rVert}
\newcommand{\snorm}[1]{\lVert#1\rVert}
\newcommand{\bvec}[1]{\boldsymbol{\mathsf{#1}}}
\newcommand{\bcov}[1]{\boldsymbol{#1}}
\newcommand{\bdua}[1]{\boldsymbol{\check{#1}}}
%\newcommand{\bdov}[1]{\boldsymbol{\breve{#1}}}
\newcommand{\bdov}[1]{\breve{#1}}
%\newcommand{\bten}[1]{\boldsymbol{\mathfrak{#1}}}
\newcommand{\bten}[1]{\boldsymbol{\mathfrak{#1}}}
\newcommand{\forany}{\tilde{\forall}}
\newcommand{\qed}{$\overset{\circ}{.}\;$}
\newcommand{\bigo}{\mathcal{O}}
\newcommand{\spn}{\mathrm{spn}\,}
\newcommand{\krn}{\mathrm{krn}\,}
\newcommand{\biy}{\leftrightarrow}
\newcommand{\ugh}{\mathbin{\textlinb{\BPwheel}}}

\newcommand\bigeye{\ensurestackMath{\stackinset{c}{}{c}{-.3pt}%
  {\bullet}{\scriptstyle\bigcirc}}}
\newcommand\eye{\scalerel*{\bigeye}{x}}
%\newcommand*{\fisheye}{%
%    \mathbin{%
%        \ooalign{$\circledcirc$\cr\hidewidth$\bullet$\hidewidth}%
%    }%
%}

% double angle-bracket
% https://tex.stackexchange.com/questions/79657/how-to-get-double-angle-bracket-without-using-mnsymbol-package
\makeatletter
\newsavebox{\@brx}
\newcommand{\llangle}[1][]{\savebox{\@brx}{\(\m@th{#1\langle}\)}%
  \mathopen{\copy\@brx\mkern2mu\kern-0.9\wd\@brx\usebox{\@brx}}}
\newcommand{\rrangle}[1][]{\savebox{\@brx}{\(\m@th{#1\rangle}\)}%
  \mathclose{\copy\@brx\mkern2mu\kern-0.9\wd\@brx\usebox{\@brx}}}
\makeatother

% https://tex.stackexchange.com/questions/297362/bar-through-letter-in-mathmode
%\usepackage{calc}
\newsavebox\CBox
\newcommand\hcancel[2][0.5pt]{%
  \ifmmode\sbox\CBox{$#2$}\else\sbox\CBox{#2}\fi%
  \makebox[0pt][l]{\usebox\CBox}%  
  \rule[0.77\ht\CBox-#1/2]{\wd\CBox}{#1}}
\newcommand{\cb}{\hcancel{b}}
\newcommand{\cd}{\hcancel{d}}
\newcommand{\ct}{\hcancel{\partial}}
%\newcommand{\supp}{\mathsf{supp}}
\newcommand{\supp}{\mathrm{supp}}
\newcommand{\csupp}{\overline{\mathrm{supp}}}

\begin{document}

% Use the \preprint command to place your local institutional report
% number in the upper righthand corner of the title page in preprint mode.
% Multiple \preprint commands are allowed.
% Use the 'preprintnumbers' class option to override journal defaults
% to display numbers if necessary
%\preprint{}

%Title of paper
\title{
M\'etodos Matem\'aticos de la F\'isica II
}

% repeat the \author .. \affiliation  etc. as needed
% \email, \thanks, \homepage, \altaffiliation all apply to the current
% author. Explanatory text should go in the []'s, actual e-mail
% address or url should go in the {}'s for \email and \homepage.
% Please use the appropriate macro foreach each type of information

% \affiliation command applies to all authors since the last
% \affiliation command. The \affiliation command should follow the
% other information
% \affiliation can be followed by \email, \homepage, \thanks as well.
\author{Juan I. Perotti}
\email[]{juan.perotti@unc.edu.ar}
%\homepage[]{Your web page}
%\thanks{}
%\altaffiliation{}
%\affiliation{}
\affiliation{Instituto de F\'isica Enrique Gaviola (IFEG-CONICET), Ciudad Universitaria, 5000 C\'ordoba, Argentina}
\affiliation{Facultad de Matem\'atica, Astronom\'ia, F\'isica y Computaci\'on, Universidad Nacional de C\'ordoba, Ciudad Universitaria, 5000 C\'ordoba, Argentina}

% \author{Colaborador}
% \affiliation{Instituto de F\'isica Enrique Gaviola (IFEG-CONICET), Ciudad Universitaria, 5000 C\'ordoba, Argentina}
% \affiliation{Facultad de Matem\'atica, Astronom\'ia, F\'isica y Computaci\'on, Universidad Nacional de C\'ordoba, Ciudad Universitaria, 5000 C\'ordoba, Argentina}
% \email[E-mail: ]{xxx.yyy@unc.edu.ar }

%Collaboration name if desired (requires use of superscriptaddress
%option in \documentclass). \noaffiliation is required (may also be
%used with the \author command).
%\collaboration can be followed by \email, \homepage, \thanks as well.
%\collaboration{Claudio J. Tessone}
%\noaffiliation

\date{\today}

\begin{abstract}
Bla bla bla... resumen...
\end{abstract}

% insert suggested keywords - APS authors don't need to do this
%\keywords{}

%\maketitle must follow title, authors, abstract, and keywords
\maketitle
\onecolumngrid

\section{Metric space}

A set $X$ is called a {\em metric space} if it is equiped with a notion of {\em metric distance}, i.e. a function $d:X\times X\to \RR$ such that, for any three {\em points} $x$, $y$ and $z$ in $X$, it holds:
\begin{enumerate}
\item $d(x,y) \geq 0$, i.e. it is {\em non-negative}.
\item $d(x,y) = 0$ implies $x=y$, i.e. it is {\em positive definite}.
\item $d(x,y) = d(y,x)$, i.e. it is {\em symmetric}.
\item $d(x,z) \leq  d(x,y) + d(y,z)$, i.e. it satisfies the {\em triangular inequality}.
\end{enumerate}

{\em Example:} The set $\RR^n$ with the euclidean distance $d(x,y) = \sqrt{\sum_{i=1}^n (x_i-y_i)^2}$ is a metric space.

\subsection{Open ball}

The set
$$
B_r(x)
=
\{y\in \RR^n: d(x,y)<r\}
$$
is called the {\em open ball} of $\RR^n$ centered at $x$ and of radious $r>0$.

\subsection{Open set}

A subset $U$ of $\RR^n$ is {\em open} if for every $x\in U$ there exists an open ball $B_r(x) \subset U$.

\subsection{Complement set}

Consider two subsets $X$ and $Y$ of $\RR^n$ such that $X\subset Y$.
The {\em complement} $X\div Y$ of $X$ in $Y$ is the subset of $\RR^n$ given by
$$
X\backslash Y := \{y\in Y:y\notin X\}.
$$
In particular, we can think of the complement $X\backslash \RR^n$.

\subsection{Closed set}

A subset $F$ of $\RR^n$ is {\em closed} if its complement $F\backslash \RR^n$ is open.

\subsection{Limiting or accumulation point}

A point $x$ is a {\em limiting point} of a subset $A$ of $\RR^n$ if, for every $r>0$, it holds that
$$
(B_r(x)\backslash \{x\})\cap A \neq \emptyset
$$
or, in other words, there are points in $A$ that are arbitrarily close to $x$.

{\em Example:}
The points $0$ and $1$ are limiting points of the open interval $(0,1)$.

\subsection{Closure set}

The {\em closure} $\overline{A}$ of a subset $A$ of $\RR^n$ is defined by
$$
\overline{A}
=
A \cup \{\mbox{all limiting points of $A$}\}.
$$

{\em Examples:} the closed interval $[0,1]$ is the closure of itself and the subsets $(0,1)$, $[0,1)$, $(0,1]$ and $(0,1/2)\cup (1/2,1)$.

\section{Topological space}

Previous definitions of open and closed sets depend on the existence of a metric distance.
It is possible to generalize the formalism using a weaker, i.e. more general, requirement.
Namely, the notions of open and closed sets can be defined when the set is equipped with a notion of topology, which provides a weaker notion of space, than that of a metric.

%\subsection{Topological space}

A {\em topological space} is a set $X$ equipped with a collection $T$ of subsets of $X$ called {\em open sets} that satisfy the following properties:
\begin{enumerate}
\item
$X$ and $\emptyset$ belong to $T$.
\item
Arbitrary unions of open sets are open.
Formally, for any $R\subset T$, then 
$$
\cup_{U\in R}\, U \in T.
$$
\item
Finite intersections of open sets are open.
Formally, for any $R\subset T$ such that $|R|\in \NN$, then
$$
\cap_{U\in R}\, U \in T.
$$
\end{enumerate}
In other words, the notion of openness is a primite definition.

The notions of complement set and closed set are analogous to those previously defined.
Namely, the complement of a subset $S$ of $X$ is the subset $S\backslash X := \{x\in X:x\notin S\}$ of $X$, and a subset $Y$ of $X$ is closed if its complement $Y\backslash X$ is open.

\subsection{Closure set}

The closure $\overline{A}$ of a subset $A$ of $X$ is defined as
$$
\overline{A}
=
\cap_{F:A\subset F,\,F\, \text{is closed.}}\,F
$$

\section{Test functions}

\subsection{The closure of a set}

The {\em closure} of a subset $S$ of $\RR^n$ is the set
$$
\overline{S}
=
...
$$

\subsection{The support of a function}

The {\em support} of a function $\varphi : \RR^n\to \RR$, is the subset $K$ of $\RR^n$ given by:
$$
\supp(\varphi)
=
\left\{
x\in \RR^n
:
\varphi(x) \neq 0
\right\}.
$$
%While the {\em clossed support} is

\subsection{The space of smooth compactly supported functions}

Let $C_c^{\infty}(\RR^n)$ denote the {\em space of smooth compactly supported functions}.
That is, if $\varphi\in C_c^{\infty}(\RR^n)$, then:
\begin{enumerate}
\item
$\varphi : \RR^n \to \RR$,

\item
$\varphi \in C^{\infty}(\RR^n)$, i.e. it is infinitely differentiable.

\item

\end{enumerate}


\end{document}
%
% ****** End of file apstemplate.tex ******


